%%%%

%%%   Самарский государственный технический университет

%%%   Кафедра прикладной математики и информатики

%%%

%%%   Преамбула для статьи в журнал "Вестник Самарского государственного

%%%   технического университета. Серия: «Физико-математические науки»"

%%%   Ориентирована под MikTEx 2.4 for Win32

%%%

%%%   Хотя сам журнал собирается под GNU/Linux на дистрибутиве TeXLive



\documentclass[11pt,a4paper,twoside]{article}



\usepackage[cp1251]{inputenc}

\usepackage[T2A]{fontenc}

\usepackage[english,russian]{babel}

\usepackage{samgtubib}

\usepackage[dvips]{graphicx}

\usepackage[multiple]{footmisc}



\usepackage{samgtu}

\renewcommand{\VestnikYear}{2009} % Год издания Вестника

\renewcommand{\VestnikNumber}{2\,(19)} % Номер Вестника

\renewcommand{\VestnikMonth}{Ноябрь} % Месяц выхода Вестника

\renewcommand{\VestnikData}{01.11.09} % Дата подписания Вестника в печать

\renewcommand{\VestnikUPL}{26,64} % Усл. печ. л.

\renewcommand{\VestnikUIL}{25,73} % Уч.-изд. л.

\renewcommand{\VestnikRegNumber}{479/09} % Регистрационный номер

\renewcommand{\VestnikOrder}{994} % Номер заказа







%

%\begin{document}





\renewcommand{\firstpage}{31}

\renewcommand{\lastpage}{38}



\udk{517.965:514.74}

\msc{34K05}



\title{Об одном классе функционально-дифференциальных уравнений}

{Об одном классе функционально-дифференциальных уравнений}



\author{В.\,А.~Кыров}

{К\,ы\,р\,о\,в~В.\,А.}



\institute{Горно-Алтайский государственный университет, \\

649000, Россия, Горно-Алтайск, ул. Ленкина, 1.}





\email{E-mail}{kfizika@gasu.ru}



\otherinf{\fullauthor{Владимир Александрович Кыров}\regalia{к.ф.-м.н., доц.} \position{доцент} \dept{каф. физики и~методики  преподавания физики}}



\keywords{функционально-дифференциальное уравнение, метрическая функция, феноменологически симметричная геометрия, геометрия

Гельмгольца}



\workabstract{В работе рассматриваются специальные функционально"=дифференциальные уравнения, возникающие в~геометрии, для метрической функции. Доказана теорема о виде метрической функции.  }



\engtitle{On some class of functional-differential equations}



\engauthor{V.\,A.~Kyrov}{K\,y\,r\,o\,v~V.\,A.}





\enginstitute{Gorny Altai State University,\\

1, Lenkin st.,  Gorno-Altaisk,  649000, Russia.}



\otherenginfm{\fullauthor{Vladimir A. Kyrov} \regalia{Ph.\,D. (Phys. \& Math.)} \position{Associate Professor}

\dept{Dept. of Physics and Teaching Methodology of Physics}}



\engabstract{In this paper we consider special functional-differential equations arising in geometry for the metric functions. We prove a theorem on the form of the metric functions.

}



\engkeywords{functional-differential equation, metric function, phenomenologically symmetric geometry, Helmholtz's geometry}



\date{24/III/2011}

\revisiondate{12/II/2012}



\maketitle





\Section[n]{Введение} Цель данной работы "--- расширение специальных функционально"=дифференциальных уравнений на метрическую функцию трёхмерной феноменологически симметричной геометрии:

\begin{equation}

\label{kyrov:1.6}

f(ij) = f(\theta(x_i,\,x_j,\,y_i,\,y_j),\,z_i,\,z_j),

\end{equation}

где

\begin{equation}

\label{kyrov:1.4}

\theta(ij) =  [(x_i - x_j)^2 + \varepsilon(y_i -

y_j)^2]\exp\left[2\varPhi_{\varepsilon}\left(\frac{y_i-y_j}{x_i-x_j}\right)\right],\quad \varepsilon

= + 1, -1, 0;

\end{equation}

$ \varPhi_{1}(u) = \gamma \arctg u$; 

$\varPhi_{-1}(u) = \beta\,\text{Ar(c)th}\, u$; 

$\varPhi_{0}(u) = u $; 

$\gamma$, $\beta = \rm{const}\ne 0$, $\beta \ne \pm 1$; 

$(x_i,\,y_i)$ "--- координаты точки $i\in M$;

$\theta(ij)$ "--- метрическая функция двумерной феноменологически симметричной геометрии Гельмгольца \cite{kyrov:bib1}. 

Вид метрической функции \eqref{kyrov:1.6} следует из теоремы, доказанной В.\,Х.~Львом в работе \cite{kyrov:bib2}:

{\it метрическая функция любой трехмерной феноменологически

симметричной геометрии содержит как аргумент метрическую функцию

двумерной феноменологически симметричной геометрии.}



Функционально"=дифференциальные уравнения, о которых здесь

говорится, в явном виде записываются при доказательстве теоремы.

Затем они сводятся к дифференциальным уравнениям. 

Во всех формулировках настоящей статьи предполагается существование подходящей локальной системы координат.



\smallskip



\Section[n]{Основные определения} Рассмотрим многообразие $M$,

$\dim M = m$. Пусть на $M$ задана гладкая функция $f$: $M\times

M\rightarrow \mathbb R$, называемая {\it метрической}, с~открытой и~плотной в $M\times M$ областью определения $S_f$. 

Локальные координаты в $M$ обозначим через $(x^1,\,\ldots,\,x^m)$. Пусть

выполняются следующие аксиомы \cite{kyrov:bib3}.



{\small \sc Аксиома невырожденности.} {\it 

Для любого упорядоченного набора $m+1$ точек $\langle i,\,i_1,\,\ldots,\,i_m\rangle$ из открытого и

плотного подмножества в $M^{m+1}:$ \linebreak $\langle ii_1\rangle,\ldots,$

$\langle ii_m\rangle,$ $\langle i_1i\rangle,\ldots,\langle

i_mi\rangle\in S_f$ выполняются неравенства

$$

\frac{\partial(f(ii_1),\ldots,f(ii_m))}{\partial(x^1_i,\ldots,x^m_i)}\ne0,\quad

\frac{\partial(f(i_1i),\ldots,f(i_mi))}{\partial(x^1_i,\ldots,x^m_i)}\ne0,

$$

где $(x^1_i,\ldots,x^m_i)$ "--- координаты точки $i\in M.$

}



\smallskip



{\small \sc Аксиома феноменологической симметрии.} {\it  

Для некоторой окрестности всякой последовательности точек $\langle

i_1,\ldots,i_{m+2}\rangle$ из открытого и плотного подмножества

прямого произведения $M^{m+2},$ таких что $\langle i_pi_q\rangle\in S_f$, $p,$ $q = 1,\,\ldots,\,m+2,$ $p\ne q,$ существует

достаточно гладкая функция $\varPhi:$ $\mathbb R^{(m+1)(m+2)/2}\to \mathbb R,$

$\rm{grad\,}\varPhi \ne 0,$ для которой имеет место тождество\/{\rm :}

$$

\varPhi(f(i_1i_2),\ldots,f(i_{m+1}i_{m+2})) = 0.

$$

}



\smallskip



\begin{definition}[1]Говорят, что на многообразии $M$

метрическая функция $f$ задает феноменологически симметричную

геометрию, если выполняются аксиомы \textit{невырожденности} и

\textit{феноменологической симметрии}.

\end{definition}



\smallskip



\begin{definition}[2]Гладкое локально инъективное отображение $

\lambda$: $M\rightarrow M$ называется {\it локальным движением},

если для любой пары $\langle ij\rangle\in S_f$ такой, что $\langle

\lambda(i),\,\lambda(j)\rangle\in S_f$, имеет место равенство

$$

f(\lambda(i),\,\lambda(j)) = f(ij).

$$

\end{definition}



\smallskip



Множество всех так определенных движений образует группу Ли преобразований. 

Можно показать, что размерность группы движений $m$-мерной феноменологически симметричной геометрии  равна

$m(m+1)/2$. 

Доказывается, что двухточечным инвариантом такой группы является метрическая функция, которая восстанавливается по данному инварианту \cite{kyrov:bib3}.



Алгебра Ли группы движений состоит из операторов вида

\begin{equation}

\label{kyrov:1.1}

X = X_1\partial_{x^1}+ \cdots +

X_m\partial_{x^m},

\end{equation}

причём $X_s =X_s(x^1,\,\ldots,x^m)$ "--- достаточно гладкие функции, $s = 1,\, 2, \,\ldots,\,m$, через которые записывается критерий локальной инвариантности метрической функции \cite[c.~35]{kyrov:bib4}:

\begin{equation}

\label{kyrov:1.2}

 X(i)f(ij) + X(j)(ij) = 0.

 \end{equation}

Операторы, образующие базис алгебры Ли для $m=3$, обозначим через $X$, $Y$, $Z$, $U$, $V$, $W$. 

Для геометрии Гельмгольца операторы алгебры Ли группы движений следующие \cite{kyrov:bib1,kyrov:bib5}:

\begin{equation}

\label{kyrov:1.45}

U = \partial_x,\quad V = \partial_y,\quad W = (-\varepsilon y - \alpha_{\varepsilon}x)\partial_x + (x -

\alpha_{\varepsilon}y)\partial_y.

\end{equation}



\pagebreak



\Section[N]{Основная теорема работы и её доказательство}



\begin{theorem}

Метрическая функция $f(ij)$ вида \eqref{kyrov:1.6} трёхмерной

феноменологически симметричной геометрии с точностью до локального

диффеоморфизма $\psi(f)\to f$ и в подходящих локальных координатах

имеет вид

\begin{equation}

\label{kyrov:1.9}

f(ij) = \left[(x_i - x_j)^2 + \varepsilon(y_i - y_j)^2\right]\exp\left[2\varPhi_{\varepsilon}\left(\frac{y_i-y_j}{x_i-x_j}\right) + 2z_i+2z_j\right],

\end{equation}

где $\varepsilon = +1,$ $-1,$ $0$; $ \varPhi_{1}(u) = \gamma\arctg

u;$ $\varPhi_{-1}(u) = \beta{\rm \, Ar(c)th\,} u,$ $\varPhi_{0}(u) = u;$

$\gamma,$ $\beta = {\rm const} \ne 0,$ $\beta \ne \pm1.$

\end{theorem}



\begin{proof}

Метрическую функцию будем искать в виде

\begin{multline}

\label{kyrov:3.1}

f(ij) = f\left(\bigl[(x_i - x_j)^2 + \varepsilon(y_i -

y_j)^2\bigr]\exp\left[2\varPhi_{\varepsilon}\Bigl(\frac{y_i-y_j}{x_i-x_j}\Bigr)\right],z_i,z_j\right)

= \\

=f(\theta,\,z_i,\,z_j).

\end{multline}

Она является двухточечным инвариантом шестимерной группы

преобразований, поэтому критерий инвариантности \eqref{kyrov:1.2} для базисных операторов $X$ и $Y$ алгебры Ли даёт

функционально"=дифферен\-циальные  уравнения

\begin{gather}

\label{kyrov:3.2}

2[X]e^{2\varPhi_{\varepsilon}\left(\frac{y_i-y_j}{x_i-x_j}\right)}\frac{\partial

f}{\partial \theta}+  X_{3}(i)\frac{\partial f}{\partial z_i} +

X_{3}(j)\frac{\partial f}{\partial z_j} =

0,\\

\label{kyrov:3.3}

 2[Y]e^{2\varPhi_{\varepsilon}\left(\frac{y_i-y_j}{x_i-x_j}\right)}\frac{\partial

f}{\partial \theta}+ Y_{3}(i)\frac{\partial f}{\partial z_i} +

Y_{3}(j)\frac{\partial f}{\partial z_j} =

0,

\end{gather}

где 

$$[X] {=} \bigl((x_i-x_j) - \alpha_{\varepsilon}(y_i-y_j)\bigr)\bigl(X_1(i) - X_1(j)\bigr) +

\bigl(\varepsilon(y_i-y_j) + \alpha_{\varepsilon}(x_i - x_j)\bigr)\bigl(X_2(i) - X_2(j)\bigr),$$ 

$$[Y] = \bigl((x_i-x_j) - \alpha_{\varepsilon}(y_i-y_j)\bigr)\bigl(Y_1(i) - Y_1(j)\bigr) +

\bigl(\varepsilon(y_i-y_j) + \alpha_{\varepsilon}(x_i - x_j)\bigr)\bigl(Y_2(i) - Y_2(j)\bigr).$$

Отсюда

\begin{equation}

\label{kyrov:3.35}

 \left[X\right]\exp\left({2\varPhi_{\varepsilon}\Bigl(\frac{y_i-y_j}{x_i-x_j}\Bigr)}\right) =

a(ij)\theta,\quad 

\left[Y\right]\exp \left({2\varPhi_{\varepsilon}\Bigl(\frac{y_i-y_j}{x_i-x_j}\Bigr)}\right)= b(ij)\theta,

\end{equation}

причём $a(ij)$ и $b(ij)$  не могут зависеть от $\theta$.



Заметим, что в уравнении (\ref{kyrov:3.2}) $X_3 \ne 0$, иначе

базисный оператор $X$, как  доказано в работе \cite{kyrov:bib6}, будет

являться линейной комбинацией операторов из системы \eqref{kyrov:1.4},

что недопустимо. Аналогично, $Y_3\ne 0$.



Пусть $[X] = [Y] = 0$. Тогда уравнения \eqref{kyrov:3.2} и \eqref{kyrov:3.3}

принимают вид

$$

X_{3}(i)\frac{\partial f(ij)}{\partial z_i} +

X_{3}(j)\frac{\partial f(ij)}{\partial z_j} =

0,\,Y_{3}(i)\frac{\partial f(ij)}{\partial z_i} +

Y_{3}(j)\frac{\partial f(ij)}{\partial z_j} = 0.

$$

Разрешая первое уравнение относительно ${\partial f(ij)}/{\partial z_i}$ и результат подставляя во второе уравнение, получаем

$$

(Y_{3}(i)X_{3}(j) - Y_{3}(j)X_{3}(i))\frac{\partial

f(ij)}{\partial z_j} = 0.

$$

По аксиоме невырожденности ${\partial f(ij)}/ {\partial z_j}\ne0$, поэтому $$Y_{3}(i)X_{3}(j) - Y_{3}(j)X_{3}(i) = 0.$$

Разделяя переменные, получаем $Y_{3} = aX_{3}$, где $a =

\rm{const}$. Тогда уравнения \eqref{kyrov:3.2} и \eqref{kyrov:3.3} линейно

зависимы, что недопустимо.



Пусть теперь $[X] = 0$ и $[Y]\ne0$. Тогда уравнение (\ref{kyrov:3.2})

принимает такой вид:

\begin{equation}

\label{kyrov:3.4}

X_3(i)\frac{\partial f(ij)}{\partial z_i} + X_3(j)\frac{\partial f(ij)}{\partial z_j} =

0.

\end{equation}

Здесь, как и в работе \cite{kyrov:bib6}, доказывается, что $X_3 = a(z)$. 

Тогда, обозначая  в~\eqref{kyrov:3.4} $\displaystyle  \int\frac{dz}{a(z)}\to

z$, имеем 

$$

 f(ij) = f\left(\bigl[(x_i - x_j)^2 + \varepsilon(y_i - y_j)^2\bigr]\exp\left(2\varPhi_{\varepsilon}\Bigl(\frac{y_i-y_j}{x_i-x_j}\Bigr)\right),z_i-z_j\right)

= f(\theta,w),

$$

где $w = z_i-z_j.$ Подставляя найденную функцию в (\ref{kyrov:3.3}), получаем

\begin{equation}

\label{kyrov:3.45}

[Y]\exp\left({2\varPhi_{\varepsilon}\Bigl(\frac{y_i-y_j}{x_i-x_j}\Bigr)}\right)\frac{\partial f(ij)}{\partial \theta} +

\bigl(Y_{3}(i) - Y_{3}(j) \bigr)\frac{\partial f(ij)}{\partial w} =

0.

\end{equation}

В \eqref{kyrov:3.45} $Y_{3} \ne \rm{const}$, иначе $[Y] = 0$, что недопустимо. Тогда

$$

\frac{[Y]\exp\left({2\varPhi_{\varepsilon}\left(\frac{y_i-y_j}{x_i-x_j}\right)}\right) }{Y_{3}(i)

- Y_{3}(j)} = \psi(\theta,\,w)\ne0.

$$

Из \eqref{kyrov:3.35} следует, что $\psi(\theta,\,w) = p(w)\theta$. 

Возвращаясь с найденным в~\eqref{kyrov:3.45}, получаем дифференциальное уравнение

$$

p(w)\theta\frac{\partial f(ij)}{\partial

\theta} + \frac{\partial f(ij)}{\partial w} = 0.

$$

Интегрируя это уравнение, определяем метрическую функцию

$$

 f(ij) = \theta q(w) =

[(x_i - x_j)^2 + \varepsilon(y_i -

y_j)^2]\exp\left[2\varPhi_{\varepsilon}\left(\frac{y_i-y_j}{x_i-x_j}\right)\right]q(w).

$$







Можно показать, что группа движений геометрии с найденной

метрической функцией имеет размерность меньше 6, т.\,е. геометрия не

является феноменологически симметричной.



Итак, $[X]\ne0$ и $[Y]\ne0$. Как и в работе \cite{kyrov:bib6}, можно доказать 

следующую лемму.



\begin{lemma}[1]Если в уравнении  \eqref{kyrov:3.2}

$$

[X]_{w=0}=0,

$$

то метрическая функция имеет вид $f(ij) = f(\theta,w).$

\end{lemma}



Метрические функции найденного вида рассмотрены выше и не задают феноменологически симметричную геометрию.



Итак, в уравнениях \eqref{kyrov:3.2} и \eqref{kyrov:3.3} $[X]_{w=0}\ne0$ и $[Y]_{w=0}\ne0.$



Уравнение \eqref{kyrov:3.2} можно переписать в виде

\begin{equation}

\label{kyrov:3.51}

 2[X]\exp\left({2\varPhi_{\varepsilon}\Bigl(\frac{y_i-y_j}{x_i-x_j}\Bigr)}\right) +

X_3(i)\psi_1(ij) + X_3(j)\psi_2(ij) = 0,

\end{equation}

где $\psi_1(ij) = (\partial f(ij)/\partial z_i)/(\partial f(ij)/\partial\theta)$, $\psi_2(ij) = (\partial f(ij)/\partial

z_j)/(\partial f(ij)/\partial\theta)$.



Предположим, что в \eqref{kyrov:3.51} $X_3 = c(z)$. 

Тогда с учётом \eqref{kyrov:3.45} это уравнение принимает вид

\begin{equation}

\label{kyrov:3.8}

[X]\exp\left({2\varPhi_{\varepsilon}\Bigl(\frac{y_i-y_j}{x_i-x_j}\Bigr)}\right) =

\bigl(a(z_i) + a(z_i)\bigr)\theta\ne0.

\end{equation}



\begin{lemma}[2]Решением функционального уравнения \eqref{kyrov:3.8} являются функции

$$

 X_1 = p(\alpha_{\varepsilon}y - x) +

q(\varepsilon\alpha_{\varepsilon}x -y)  + c_1,\quad  X_2 =

p(\varepsilon\alpha_{\varepsilon}x -y) - \varepsilon

q(\alpha_{\varepsilon}y - x) + c_2,

$$

где $p,$ $q,$ $c_1,$ $c_2,$ $a = \rm const.$

\end{lemma}





Подставим найденное в \eqref{kyrov:3.2} и, переобозначая $\displaystyle \int\cfrac{dz}{c(z)}\to z$, после интегрирования имеем

\begin{equation}

\label{kyrov:3.9}

f(ij) = f(\theta e^{2z_i},\,\theta e^{2z_j}) = f(u,v).

\end{equation}





Подставляя найденное в \eqref{kyrov:3.3}, получаем функциональное уравнение

\begin{equation}

\label{kyrov:3.10}

[Y+ \bar{\theta} Y_{3i}]e^{2z_i} \frac{\partial f}{\partial u} +[Y

+ \bar{\theta} Y_{3j}]e^{2z_j}\frac{\partial f}{\partial v} = 0,

\end{equation}

где $ \bar{\theta} = (x_i-x_j)^2 + \varepsilon(y_i-y_j)^2. $ 

Так как \eqref{kyrov:3.10} "--- дифференциальное уравнение, имеет место тождество

\begin{equation}

\label{kyrov:3.11}

\frac{[Y]+ \bar{\theta} Y_{3i}}{[Y]+ \bar{\theta} Y_{3j}} = \varphi(u,v).

\end{equation}





\begin{lemma}[3]В тождестве \eqref{kyrov:3.11} $\varphi(u,\,v) = -1.$

\end{lemma}



\begin{proof}

Пусть в тождестве \eqref{kyrov:3.11} $z_i = z_j = z$. 

Тогда оно принимает вид

$$

 2[Y] + \bar{\theta} (Y_{3i} + Y_{3j}) = 0.

$$

Дифференцируя по $x_i$ и $x_j$,  по $y_i$ и $y_j$, а также по $x_i$ и $y_j$, будем иметь систему

\begin{equation}

\label{kyrov:3.12}

\begin{array}{l}

-Y'_{1x_i} - Y'_{1x_j} - \alpha_{\varepsilon}Y'_{2x_i} - \alpha_{\varepsilon}Y'_{2x_j} - (Y_3(i) + Y_3(j)) + \\

\hspace{6.45cm} +(x_i - x_j)(Y'_{3x_j} - Y'_{3x_i}) = 0,

\\

\alpha_{\varepsilon}Y'_{1y_i} + \alpha_{\varepsilon}Y'_{1y_j}

-\varepsilon Y'_{2y_i} -\varepsilon Y'_{2y_j} - \varepsilon(Y_3(i)

+ Y_3(j)) + \\

\hspace{6.45cm} +\varepsilon(y_i - y_j)(Y'_{3y_j} - Y'_{3y_i}) = 0,\\

\alpha_{\varepsilon}Y'_{1x_i} - Y'_{1y_j} - \varepsilon Y'_{2x_i}

- \alpha_{\varepsilon}Y'_{2y_j} + (x_i - x_j)Y'_{3y_j} -

\varepsilon(y_i-y_j)Y'_{3x_i} = 0. \end{array}\end{equation}

Дифференцируя теперь первое уравнение системы \eqref{kyrov:3.12} по

$x_i$ и $x_j$, а также по $x_i$ и $y_j$, затем второе уравнение по

$y_i$ и $y_j$, будем иметь $ Y''_{3xx} = Y''_{3xy} = Y''_{3yy} =

0$. Интегрируя, получаем

\begin{equation}

\label{kyrov:3.13}

Y_{3} = ax +

by + c,

\end{equation}

причём коэффициенты зависят от $z$.

Подставим теперь найденное в \eqref{kyrov:3.12}:

\begin{equation*}

\begin{array}{l}

-Y'_{1x_i} - Y'_{1x_j} - \alpha_{\varepsilon}Y'_{2x_i} - \alpha_{\varepsilon}Y'_{2x_j} - (ax_i + by_i + 2c + ax_j + by_j) = 0, \\

\alpha_{\varepsilon}Y'_{1y_i} + \alpha_{\varepsilon}Y'_{1y_j}

-\varepsilon Y'_{2y_i} -\varepsilon Y'_{2y_j} - \varepsilon(ax_i +

by_i + 2c + ax_j + by_j) = 0,\\ \alpha_{\varepsilon}Y'_{1x_i} -

Y'_{1y_j} - \varepsilon Y'_{2x_i} - \alpha_{\varepsilon}Y'_{2y_j}

+ b(x_i - x_j) - \varepsilon a(y_i-y_j) = 0.

\end{array}\end{equation*}

Дифференцируя первые два уравнения по координатам, а в последнем

разделяя переменные, получаем систему уравнений, интегрируя

которую до конца, будет иметь

\begin{equation}

\label{kyrov:3.14}

\begin{array}{l}

\displaystyle

Y_1  = \frac{1}{2(\alpha_{\varepsilon}^2 + \varepsilon)}

\bigl(-(\alpha_{\varepsilon}b+\varepsilon a)(x^2 -

\varepsilon y^2) - 2\varepsilon(b+\alpha_{\varepsilon}a)xy

-\\

\hspace{4.45cm}

-2(\varepsilon c - \alpha_{\varepsilon}d)x  + 2(\varepsilon d -

\alpha_{\varepsilon}c)y + \varepsilon k - \alpha_{\varepsilon}l\bigr),

\\ 

\displaystyle

Y_2  = \frac{1}{2(\alpha_{\varepsilon}^2 +

\varepsilon)} \bigl(-(\alpha_{\varepsilon}a - b)(x^2 - \varepsilon y^2)

- 2(\varepsilon a+\alpha_{\varepsilon}b)xy-\\ 

\hspace{4.45cm}

-2(d +

\alpha_{\varepsilon}c)x + 2(c + \alpha_{\varepsilon}d)y - l +

\alpha_{\varepsilon}k\bigr),

\end{array}

\end{equation}

причём коэффициенты зависят от $z$. 

Умножая результат на исходное выражение и учитывая формулы \eqref{kyrov:3.13} и \eqref{kyrov:3.14}, получаем $a = b = c = d = k = l = \text{const}$. Затем, подставляя \eqref{kyrov:3.13} и \eqref{kyrov:3.14}  в \eqref{kyrov:3.11}, получаем $\varphi(u,v) = -1$. Лемма~3 доказана. \end{proof}



\smallskip



Тогда, подставляя \eqref{kyrov:3.13} и \eqref{kyrov:3.14}   в \eqref{kyrov:3.10} и интегрируя, получаем метрическую функцию \eqref{kyrov:1.9}. 

Базисные операторы алгебры Ли группы геометрии с~метрической функцией \eqref{kyrov:1.9} фактически найдены при доказательстве этой леммы. Для

их явной записи в формулах (\ref{kyrov:3.14}) и (\ref{kyrov:3.13})

коэффициентам необходимо придать значения 0 и 1.



Итак, в уравнении (\ref{kyrov:3.51}) имеем $\bigl[X_3(i) - X_3(j)\bigr]_{z_i=z_j=z}\ne0$. 

Если его переписать для пары $\langle ji\rangle$, то, как и выше, получаем  $\psi_1(ij) - \psi_2(ji) = 0$.

Обозначим $\psi_1(ij) = \psi(ij)$, $\psi_2(ij) = \psi(ji)$. 

Тогда уравнение (\ref{kyrov:3.51}) переписывается так:

\begin{equation}

\label{kyrov:3.135}

2[X]\exp\left({2\varPhi_{\varepsilon}\Bigl(\frac{y_i-y_j}{x_i-x_j}\Bigr)}\right) + X_3(i)\psi(ij) + X_3(j)\psi(ji) = 0.

\end{equation}



\begin{lemma}[4]В уравнении \eqref{kyrov:3.135}

$$

\psi\bigr|_{w=0} = c(z)\theta\ne0,\quad X_3 = p(z)x + q(z) y + d(z).

$$

\end{lemma}



\begin{proof}

Введём  в (\ref{kyrov:3.135}) подстановку $z=z_i=z_j$.

Обозначим $X_{1,2,3}\bigr|_{w=0}(i) = X_{1,2,3}(x_i,y_i,z)$, $\psi\bigr|_{w=0}(ij) = \psi(\theta,z,z)$. 

Заметим, что $\bar{\psi}(ij) = \bar{\psi}(ji)$. 

Тогда из \eqref{kyrov:3.35} следует, что $\psi\bigr|_{w=0} = c(z)\theta$. 

Подставляя найденное в уравнение \eqref{kyrov:3.135} при $w=0$ и~учитывая результат предыдущей леммы, получаем 

$X_3 = p(z)x + q(z) y + d(z).$

Лемма 4 доказана.

\end{proof}



\smallskip



Как известно, базис алгебры Ли группы движений трёхмерной феноменологически симметричной геометрии состоит из шести

операторов, три из которых "--- $U$, $V$, $W$ "--- известны, а три "--- $X$, $Y$, $Z$ "--- надо найти. 

Оператор $X$ удовлетворяет уравнению \eqref{kyrov:3.135}, а $Y$ и $Z$ "--- уравнениям 

\begin{gather}

\label{kyrov:2.231}

 2[Y]\exp \left({2\varPhi_{\varepsilon}\Bigl(\frac{y_i-y_j}{x_i-x_j}\Bigr)}\right) + Y_3(i)\psi(ij) + Y_3(j)\psi(ji) = 0, \\

\label{kyrov:2.241}

2[Z]\exp\left({2\varPhi_{\varepsilon}\Bigl(\frac{y_i-y_j}{x_i-x_j}\Bigr)}\right)+ Z_3(i)\psi(ij) + Z_3(j)\psi(ji) = 0,

\end{gather}

причём

$\bigl[Y_3(i) - Y_3(j)\bigr]_{z_i=z_j=z}\ne0$, $\bigl[Z_3(i) - Z_3(j)\bigr]_{z_i=z_j=z}\ne0$, иначе приходим к~рассмотренному ранее  случаю.



Итак, для \eqref{kyrov:3.135} имеем

$$

[X]_{w=0}\exp \left({2\varPhi_{\varepsilon}\Bigl(\frac{y_i-y_j}{x_i-x_j}\Bigr)}\right) + (p_1(x_i+x_j) + q_1(y_i+y_j) + d_1)\bar{\theta} = 0.

$$

Аналогичный результат получаем для уравнений \eqref{kyrov:2.231} и \eqref{kyrov:2.241}:

\begin{gather*}

[Y]_{w=0}\exp\left({2\varPhi_{\varepsilon}\Bigl(\frac{y_i-y_j}{x_i-x_j}\Bigr)}\right) + (p_2(x_i+x_j) + q_2(y_i+y_j) + d_2)\bar{\theta} = 0,\\

[Z]_{w=0}\exp\left({2\varPhi_{\varepsilon}\Bigl(\frac{y_i-y_j}{x_i-x_j}\Bigr)}\right) + (p_3(x_i+x_j) + q_3(y_i+y_j) + d_3)\bar{\theta} = 0.

\end{gather*}

Комбинируя эти уравнения, приходим к ранее рассмотренному случаю

(лемма 4), который приводит к метрической функции \eqref{kyrov:1.9}.

Теорема доказана полностью.

\end{proof}



\thanks{Выражаю благодарность профессору  Геннадию

Григорьевичу  Михайличенко за поддержку темы исследования.}





\begin{thebibliography}{99}



\RBibitem{kyrov:bib1}

\by Кыров~В.\,А.

\paper Гельмгольцевы пространства размерности два

\jour Сиб. матем. журн.

\yr 2005

\vol 46

\issue 6

\pages 1341--1359

\transl англ. пер.:~

\paper Two-Dimensional Helmholtz Spaces

\by  Kyrov~V.\,A.

\jour Siberian Math. J.

\yr 2005

\vol 46

\issue 6

\pages 1082--1096



\RBibitem{kyrov:bib2}

\by Лев~В.\,Х.

\paper Трёхмерные геометрии в теории физических структур

\inbook Вычислительные системы. {\rm Вып.~125}

\publaddr Новосибирск

\publ ИМ СОАН СССР

\yr 1988

\pages 90--103

\misctransl

\by Lev~V.\,Kh.

\paper Three-Dimensional Geometries in Physical Structures Theory

\inbook Vychislitel'nye Sistemy. {\rm Issue~125}

\publaddr Novosibirsk

\publ IM SOAN SSSR

\yr 1988

\pages 90--103









\RBibitem{kyrov:bib3}

\by Михайличенко~Г.\,Г.

\paper О групповой и феноменологической симметриях в геометрии

\jour Докл. АН СССР

\yr 1983

\vol 269

\issue 2

\pages 284--288

\transl англ. пер.:~

\by Mikhailichenko~G.\,G.

\paper  On Group and Phenomenological Symmetries in Geometry 

\jour Dokl. Soviet. Math

\yr 1983

\vol 27

\issue 2

\pages 325--329





\RBibitem{kyrov:bib4}

\by Овсянников~Л.\,В.

\book Групповой анализ дифференциальных

уравнений

\publaddr M.

\publ Наука

\yr 1978

\totalpages 399

\misctransl

\by Ovsjannikov~L.\,V. 

\book Group analysis of differential equations

\publ Nauka

\publaddr Moscow

\yr 1978

\totalpages 399



\RBibitem{kyrov:bib5}

\by Кыров~В.\,А. \nofrills

\paperinfo  Шестимерные алгебры Ли групп движений трехмерных феноменологически симметричных геометрий: приложение к книге \nofrills

\byy Г.~Г.~Михайличенко \nofrills

\inbook  ``Полиметрические геометрии'' \nofrills.

\publaddr Новосибирск

\publ Новосиб. гос. ун-т

\yr 2001

\pages 116--143

\misctransl 

\by Kyrov~V.\,A. \nofrills

\paperinfo Six-dimensional Lie algebras  of  movements groups three-dimensional  phe\-no\-meno\-logically symmetric geometries. Appendix to the book \nofrills

\inbook ``Polymetric Geometries''\nofrills.

\byy G.~G.~Mikhailichenko \nofrills

\publaddr Novosibirsk

\publ Novosib. Gos. Un-t

\yr 2001

\pages 116--143





\RBibitem{kyrov:bib6}

\by Кыров~В.\,А.

\paper Функциональные уравнения в~псевдоевклидовой геометрии

\jour Сиб. журн. индустр. матем.

\yr 2010

\vol 13

\issue 4

\pages 38--51

\misctransl

\by Kyrov~V.\,A.

\paper Functional equations in pseudo-Euclidean geometry

\jour Sib. Zh. Ind. Mat.

\yr 2010

\vol 13

\issue 4

\pages 38--51





\end{thebibliography}



\noindent

\hskip6.5cm \thedate \par\noindent

\hskip6.5cm\therevdate







\makeengtitlem





\newpage







%

%\end{document}

